\scanpage{208}

Considérons la suite exacte \unsure{(1.2)|(2.2)}, et $u^*(P^c_{X/S})\to P^c_{X/S}$, je dis qu'en fait

\[
u^*(P^c_{X/S})\longrightarrow P^{ci}_{X/S},\tag*{(2.3)}
\]

\deleted{En effet} dès que l'on a

\[
f_*^{\mathrm{fl}}(\Omega^1_{X/S})=f_*^{\mathrm{fl}}(\underline Z^1_{X/S})\tag*{(2.4)}
\]

\margin{$2$}

(Attention, \deleted{dans la \unsure{définition|différentielles}} \inserted{la notation} de $f_*^{\mathrm{fl}}(\underline Z^1_{X/S})$ \deleted{est faisceau} \inserted{un abus de} $\underline Z^1_{\unsure{X/S|T/S}}$ lang.\ \unsure{pr|les} pour $\mathrm{Ker}\,(f_*^{\mathrm{fl}}(\underline\Omega^1_{X/S})\longrightarrow f_*^{\mathrm{fl}}(\underline\Omega^2_{X/S}))$], et que

\deleted{En effet, en \unsure{lisant} $u^*(P^c)\longrightarrow$}

\[
\deleted{f_*^{\mathrm{fl}}(\unsure{\Omega})} V=f_*(\underline\Omega^2_{X/S})\ \deleted{\unsure{\text{est}}}\ \deleted{\text{est}}\ \text{coreprésentable par un module}\tag*{(2.5)}
\]
quasi-cohérent, \uline{ou} représentable par un fibré vectoriel.

En effet, on a alors $\mathrm{Hom}(B,V\deleted{f_*(\Omega^{\unsure{1}}_{X/S})})=0$, donc $\mathrm{Hom}(u^*(P^c),V)\longrightarrow\mathrm{Hom}(f_*^{\mathrm{fl}}(\Omega^1_{X/S})/{-},V)$ est injectif, \deleted{donc} or l'hom obtenu est la différentielle ext. $f_*^{\mathrm{fl}}(\Omega^1_{\unsure{Y|X}/S})\longrightarrow f_*^{\mathrm{fl}}(\Omega^2_{X/S})$, qui est supposée nulle \ldots.

En particulier, sous les conditions favorables (2.2), et moyennant (2.4) et (2.5), on a

\[
(P^{c\unsure{0|o}})\subset P^{ci}\qquad\text{i.e.}\qquad P^{c0}=P^{ci0}\tag*{(2.6)}
\]

\textcircled{3}. Supposons que $p\cdot1_S=0$ ($p$ un nb premier). On définit alors

\[
P^{cip}_{X/S}\subset P^{ci}_{X/S}
\]

le sous-faisceau correspondant aux modules à connexions \uline{compatibles aux puiss.\ $p^{\mathrm{ièmes}}$} \deleted{à puissances $p^{\mathrm{ièmes}}$ nulles}. Supposons que $X/S$ ait loc. une $p$-base, alors on a

\[
0\longrightarrow P^{cip}_{X/S}\longrightarrow P^{ci}_{X/S}
\xrightarrow[\substack{\text{tenseur de}\\\deleted{\unsure{\text{défaut de}}}\\\text{$p$-courbure}\\=\deleted{\unsure{\text{obstruction}}}\,\text{de comp. aux puiss. $p^{\mathrm{ièmes}}$}}]{\overset{\inserted{\text{$p$-courbure}}}{K}}
 f_*^{\mathrm{fl}}(\Omega^{1\otimes p}_{X/S})\tag*{(3.1)}
\]

% archive sheet number: 207
